\par\Needspace{4\baselineskip}
\begin{center}
\hypertarget{native-059-article041}{}\pdfbookmark[2]{41.}{article-041}41.
\end{center}
It remains to give a more convenient algorithm for the inverse problem as well, namely to determine the time from the true anomaly. For this purpose we have added a new column for $T$ to our table. First compute $T$ from $v$ by the formula $T=\frac{1-e}{1+e}\operatorname{tang}\tfrac{1}{2}v^2$. With $T$ as argument, take $A$ and $\log B$ from our table, or, more precisely and even more conveniently, take $C$ and $\log B$ and obtain $A$ by the formula $A=\frac{(1+C)T}{1+\frac{4}{5}T}$. Finally, from $A$ and $B$ obtain $t$ by formula 1) of Art. 37. If one wishes to use \textsc{Barker}'s table here as well, it saves less work in this inverse problem. The following is then obtained directly, independently of $A$
\[
\operatorname{tang}\tfrac{1}{2}w=\operatorname{tang}\tfrac{1}{2}v\,\mathord{.}\,\BareRoot{\frac{1+C}{\gamma\gamma(1+\frac{4}{5}T)}},
\]
and hence the time $t$ is obtained by multiplying the mean motion corresponding to the true anomaly $w$ in \textsc{Barker}'s table by $\frac{B}{\alpha}\cdot$
